% ==========================================
% PHAN B - RIPV2 TRONG AS A
% ==========================================

\section{Phần B -- Cấu hình RIPv2 trong AS A}

\subsection{Mục tiêu}

Phần này cấu hình RIPv2 trên A1, A2 và A3
để trao đổi các tuyến nội bộ trong AS A,
bao gồm LAN-A và các địa chỉ Loopback0.
A2 và A3 đều quảng bá default route vào RIP
để A1 có đường đi tới các mạng bên ngoài.
Các bản tin cập nhật RIP chỉ được gửi
trên ba liên kết nội bộ L1--L3.

\subsection{Triển khai cấu hình}

\subsubsection{Ý nghĩa các lệnh chính}

Trong Cisco IOS, lệnh \texttt{network} của RIP
sử dụng địa chỉ mạng theo lớp (classful).
Lệnh này bật RIP trên các interface có
địa chỉ thuộc mạng được khai báo.
Ở bài này, các địa chỉ interface đều thuộc
mạng lớp A \texttt{10.0.0.0/8}, nên nhóm dùng:

\begin{lstlisting}
network 10.0.0.0
\end{lstlisting}

Việc khai báo theo mạng classful không có nghĩa
các tuyến được gửi đi đều mang mask /8.
Nhóm dùng \texttt{version 2} để bản tin RIP
mang thông tin subnet mask, đồng thời dùng
\texttt{no auto-summary} để tắt cơ chế
tự động tổng hợp tuyến theo ranh giới lớp mạng.
Nhờ đó, các subnet /30, /32 và /26
được trao đổi với đúng subnet mask.

Do các cổng nối sang AS khác cũng có địa chỉ
thuộc \texttt{10.0.0.0/8}, cần giới hạn
những interface được gửi cập nhật RIP:

\begin{itemize}
    \item \texttt{passive-interface default}:
    Không gửi bản tin cập nhật RIP
    ra các interface theo mặc định.

    \item \texttt{no passive-interface <interface>}:
    Cho phép gửi lại cập nhật RIP
    qua interface được chỉ định.
    Nhóm chỉ mở FastEthernet0/0 và
    FastEthernet0/1 trên A1, A2, A3,
    tương ứng với các liên kết L1--L3.

    \item \texttt{default-information originate}:
    Phát sinh và quảng bá default route
    \texttt{0.0.0.0/0} vào RIP.
    Lệnh này được cấu hình trên cả A2 và A3.
\end{itemize}

Passive interface chỉ kiểm soát việc gửi
cập nhật trên interface đó, không tự loại
mạng của interface khỏi các bản tin gửi
qua interface khác. Vì vậy, Loopback0
và LAN-A vẫn được quảng bá dù không
gửi bản tin RIP ra các interface Loopback.
Tương tự, các subnet của L4 và L5
vẫn có thể xuất hiện trong bảng RIP nội bộ,
nhưng RIP không được gửi sang ISP qua hai cổng này.

\subsubsection{Cấu hình trên A1}

A1 trao đổi cập nhật RIP với A2 qua L1
và A3 qua L2. Cấu hình như sau:

\begin{lstlisting}
conf t
router rip
 version 2
 no auto-summary
 network 10.0.0.0
 passive-interface default
 no passive-interface FastEthernet0/0
 no passive-interface FastEthernet0/1
exit
end
write memory
\end{lstlisting}

LAN-A \texttt{10.171.126.0/26} và Loopback0
\texttt{10.171.26.1/32} của A1 nằm trong
phạm vi khai báo \texttt{network 10.0.0.0},
nên được đưa vào các bản tin cập nhật
gửi tới A2 và A3.

\subsubsection{Cấu hình trên A2 và A3}

A2 và A3 dùng cùng một nhóm lệnh RIP dưới đây.
Ngoài việc trao đổi tuyến nội bộ,
cả hai router đều có thêm lệnh
\texttt{default-information originate}:

\begin{lstlisting}
conf t
router rip
 version 2
 no auto-summary
 network 10.0.0.0
 passive-interface default
 no passive-interface FastEthernet0/0
 no passive-interface FastEthernet0/1
 default-information originate
exit
end
write memory
\end{lstlisting}

Trên A2, hai cổng được mở RIP nối tới
A1 qua L1 và A3 qua L3.
Trên A3, chúng nối tới A1 qua L2
và A2 qua L3.
Cổng FastEthernet1/0 nối sang ISP
vẫn là passive trên cả hai router.

Default route giúp A1 chuyển gói tin
không có tuyến cụ thể tới router biên.
Đường đi tiếp ra ISP được A2 và A3
lựa chọn bằng BGP ở các phần sau.

\subsection{Kiểm chứng kết quả}

\subsubsection{Bảng định tuyến RIP trên A1}

\begin{lstlisting}
show ip route rip
\end{lstlisting}

\partzeroimage
{images/configuration/part_B/A1_route_rip.png}
{Các tuyến học qua RIPv2 trên A1-7126}
{fig:partB-a1-route}

Hình \ref{fig:partB-a1-route} cho thấy A1 học
được Loopback0 của A2 là \texttt{10.171.26.2/32}
qua next-hop \texttt{10.71.26.2},
và Loopback0 của A3 là \texttt{10.171.26.3/32}
qua next-hop \texttt{10.71.26.6}.
Giá trị \texttt{[120/1]} biểu thị AD của RIP
là 120 và metric tới các mạng này là 1 hop.

A1 cũng có tuyến \texttt{R* 0.0.0.0/0}
qua A2, với next-hop \texttt{10.71.26.2}.
Đây là default route được ghi nhận
tại thời điểm chụp; ảnh không thể hiện
đồng thời default route qua A3.
LAN-A và Loopback0 của chính A1
là các mạng kết nối trực tiếp,
nên không xuất hiện trong output
\texttt{show ip route rip}.

\subsubsection{Thông tin giao thức trên A2}

\begin{lstlisting}
show ip protocols
\end{lstlisting}

\partzeroimage
{images/configuration/part_B/A2_protocols.png}
{Thông tin giao thức RIP trên A2-7126}
{fig:partB-a2-protocols}

Hình \ref{fig:partB-a2-protocols} xác nhận
A2 gửi và nhận RIP phiên bản 2,
không tự động tổng hợp mạng,
và dùng khai báo \texttt{network 10.0.0.0}.
Hai cổng FastEthernet0/0 và FastEthernet0/1
tham gia trao đổi cập nhật;
FastEthernet1/0 nằm trong danh sách passive,
nên không gửi RIP sang AS B.
Các nguồn thông tin định tuyến
\texttt{10.71.26.1} và \texttt{10.71.26.10}
lần lượt tương ứng với A1 và A3.

\newpage
